\section{Introduction}
\label{sec:introduction}

Modern enterprise networks rely almost universally on Microsoft Active Directory (AD) to administer centralized identity, policy enforcement, and authentication services. Within this ecosystem, Active Directory Certificate Services (ADCS) operates as the enterprise Public Key Infrastructure (PKI), providing automated enrollment, cryptographic identity certificates, and Kerberos public-key pre-authentication (PKINIT, RFC 4556~\cite{rfc4556}) for millions of principals. Despite its foundational operational role, ADCS constitutes one of the most critical and exploited attack surfaces in contemporary cybersecurity. Research formalized by SpecterOps in 2021~\cite{specterops2021certified} identified a comprehensive taxonomy of certified pre-owned escalation vectors (ESC1 through ESC13), wherein subtle misconfigurations in certificate templates, access control lists (ACLs), and Certificate Authority (CA) registration flags allow low-privileged domain principals to compromise entire domain forests.

\subsection{The Multi-Faceted Vulnerability Landscape of ADCS}
In an ADCS deployment, certificate issuance is governed by Active Directory objects known as \textit{Certificate Templates} (stored in the LDAP configuration naming context). When an enterprise template exposes an insecure combination of attributes—such as the \texttt{CT\_FLAG\_ENROLLEE\_SUPPLIES\_SUBJECT} bitflag alongside an Extended Key Usage (EKU) permitting Client Authentication (\texttt{1.3.6.1.5.5.7.3.2}) or Smart Card Logon (\texttt{1.3.6.1.5.2.3.4}), without requiring CA administrator approval—any domain principal granted enrollment permissions can request an X.509 certificate specifying an arbitrary Subject Alternative Name (SAN). The attacker can supply the user principal name (UPN) of a Domain Administrator (e.g., \texttt{Administrator@domain.local}), obtain a cryptographically valid certificate signed by the enterprise root CA, and submit that certificate to the Kerberos Key Distribution Center (KDC) via PKINIT to request a Ticket Granting Ticket (TGT), achieving instant domain takeover.

While static privilege audit tools such as BloodHound~\cite{robbins2017bloodhound} and Certipy~\cite{certipy2021} have gained widespread industrial adoption, defending enterprise ADCS infrastructure against sophisticated adversaries reveals three structural failure modes:

\begin{enumerate}
    \item \textbf{Inductive Shortcut Memorization in Static GNNs:} Recent academic efforts have applied Graph Neural Networks (GNNs) to automate vulnerability discovery over identity multigraphs~\cite{guo2022practical,vasquez2024identity}. However, we discover that standard GNN formulations (GCN, GAT, and additive-skip architectures) succumb to severe \textit{inductive shortcut collapse}. Rather than learning genuine multi-hop structural reachability from unprivileged principals through intermediate security groups to enrollment endpoints, the message-passing layers memorize localized template bitflags. When evaluated against adversarial \textit{hard negatives} (configurations that possess dangerous template flags but are structurally unreachable due to enrollment restrictions or issuance requirements), standard GNNs collapse to an alarming $1.59\%$ accuracy ($1/63$), generating false-alarm cascades that paralyze Security Operations Centers (SOCs).
    
    \item \textbf{Static Snapshot Blindness to Temporal Mutation Attacks:} In real-world enterprise environments, sophisticated adversaries rarely exploit static misconfigurations. Instead, under vectors such as ESC4 (template ACL abuse) and ESC9 (weak shadow credentials with template enrollment), an attacker with \texttt{WriteOwner} or \texttt{GenericAll} permissions dynamically mutates a benign certificate template, issues an illegitimate administrative certificate, and immediately rolls the template attributes back to their pristine baseline. We prove empirically that the pre-attack and post-rollback directory snapshots are bit-identical ($\text{SHA-256}(G_{t_0}) = \text{SHA-256}(G_{t_3})$). Consequently, \textit{any} periodic static graph scanner suffers from a $0.0\%$ post-rollback detection recall, completely missing the dormant forged certificate.
    
    \item \textbf{The CA--KDC Logging Discontinuity and SAN Concealment:} Even when SOCs enable full Windows security telemetry, current auditing architectures fail due to an unaddressed protocol discontinuity: the event schemas of Microsoft-Windows-CertificationAuthority ($\pi_{\text{CA}}$) and Kerberos-Key-Distribution-Center ($\pi_{\text{KDC}}$) are disjoint ($\pi_{\text{CA}} \cap \pi_{\text{KDC}} = \emptyset$). Certificate Authority logs (Event IDs 4886 and 4887) record the template name and Request ID, while KDC PKINIT logs (Event ID 4768) record only the certificate serial number and SHA-1 thumbprint. Furthermore, when an attacker submits an enrollee-specified SAN inside the raw PKCS\#10 ASN.1 binary CSR (\texttt{szOID\_ENROLL\_CERTTYPE\_EXTENSION}), the CA security event log suppresses the requested SAN string entirely. A defender inspecting standard SIEM logs cannot connect the unauthorized template issuance to the subsequent administrative Kerberos authentication.
\end{enumerate}

\subsection{Our Contributions}
To address these fundamental challenges, this paper presents \textbf{CertGraph}, an integrated static-to-dynamic neuro-symbolic framework that bridges topological access-control reasoning, cryptographic audit provenance, and game-theoretic interdiction. The primary contributions of this work are summarized as follows:

\begin{itemize}
    \item \textbf{Diagnosis and Mathematical Cure for GNN Shortcut Collapse:} We formalize the failure of standard message passing on access control graphs and prove intrinsic attribute preservation bounds (Theorem~\ref{thm:theorem1}). We introduce the \textbf{Multiplicative Bilinear Reachability Gate} ($\mathbf{g}_{uv} = \sigma(\mathbf{h}_u^\top \mathbf{W}_g \mathbf{h}_v)$), which strictly modulates neighborhood aggregation by end-to-end topological reachability. This cures inductive shortcut memorization, boosting adversarial hard-negative accuracy from $1.59\%$ to \textbf{$100.00\%$ ($63/63$)} without sacrificing global detection performance.
    
    \item \textbf{Cryptographic Invariant Provenance Graph (CIPG):} We design and implement an asynchronous provenance engine that parses live Active Directory Directory Service (NTDS), CA database, and KDC telemetry. We formalize \textbf{Invariant~V} (Persistent Cryptographic Taint-Flow Automaton), which establishes a state-tracking bridge across Subject Key Identifiers (SKID), serial numbers, and SHA-1 thumbprints. On a live Windows Server 2019 cyber-range, Invariant~V achieves $100.0\%$ detection recall against split-actor collusion and 24-hour dormant certificate usage, maintains a $0.0\%$ false-positive rate across 50 administrative typo rollbacks, and reduces graph storage overhead by $69.04\%$ via meta-node collapsing at $2,228.4$~events/second.
    
    \item \textbf{Game-Theoretic Interdiction with Provable Disruption Reduction:} We formulate the \textbf{Disruption-Budgeted Stackelberg Remediation Game} between an enterprise defender and an optimal cost-minimizing attacker. We prove that Multi-Principal Access Interdiction (MPAI) is NP-complete via reduction from Directed Multiway Cut (Theorem~\ref{thm:nphardness}). Evaluated across 30 multi-forest enterprise topologies, our co-adaptive Stackelberg policy strictly dominates static greedy interdiction across tight operational budgets ($p = 3.12 \times 10^{-5}$ at $B_{\text{ops}}=6$) and achieves $100.0\%$ path neutralization at $B_{\text{ops}}=24$ while revoking an average of only \textbf{5.1 edges compared to 11.1 edges} for greedy heuristics—a \textbf{54.1\% reduction in administrative disruption}.
\end{itemize}

\subsection{Paper Organization}
The remainder of this manuscript is organized as follows: Section~\ref{sec:threat_model} formalizes the heterogeneous identity multigraph schema and demonstrates the static snapshot blind spot and ASN.1 SAN concealment on live systems. Section~\ref{sec:certgraph_arch} details the CertGraph GNN architecture and the Multiplicative Bilinear Gating cure. Section~\ref{sec:cipg_provenance} develops the CIPG provenance engine and Invariant~V automaton. Section~\ref{sec:game_theoretic} details the game-theoretic Stackelberg defense and NP-hardness proof. Section~\ref{sec:evaluation} presents empirical results across synthetic and live GOAD cyber-range testbeds. Section~\ref{sec:related_conclusion} discusses related literature and concludes the paper.
